\section{Resumen del Proyecto}
En este trabajo se estudia de forma numérica la inestabilidad de Kelvin-Helmholtz (KHI por su nombre en inglés \textit{Kelvin-Helmholtz Instability}) en dos dimensiones, utilizando el marco de la Magnetohidrodinámica Resistiva Relativista (RRMHD). Esta inestabilidad aparece cuando hay dos flujos que se mueven en direcciones opuestas o con velocidades diferentes, y además existen contrastes marcados en densidad y presión. Esto genera una interfaz en donde se pueden desarrollar estas inestabilidades.

Para estudiar este fenómeno, se trabaja con varias configuraciones en las que se cambia el valor de la resistividad, con el fin de observar cómo afecta esta propiedad al comportamiento de la inestabilidad. Las simulaciones se realizan usando el código numérico \textbf{Cueva} que resuelve las ecuaciones de la RRMHD \citep{aranguren-2013}. En \textbf{Cueva} la integración temporal se lleva a cabo mediante métodos parcialmente implícitos con el fin de tratar la rigidez de las ecuaciones de la MHD relativista resistiva (RRMHD) en el régimen ideal. En el código se implementan dos familias principales de integradores temporales: los métodos RKIMEX \citep{palenzuela-2009} y los esquemas MIRK \citep{Aloy_Cordero-Carrion:2016}. Cueva está basado en una formulación conservativa de volumen finito de las ecuaciones RRMHD, en la cual los flujos numéricos en las interfaces de las celdas se calculan utilizando ya sea el método LLF \citep{Alic2007}, el HLL \citep{HLL:1983} o el solucionador aproximado de Riemann HLLC presentado en \citep{miranda-aranguren-2018}.

Con el fin de entender mejor la relación entre resistividad e inestabilidad, se extraerán y analizarán el comportamiento de las variables globales, en toda la malla de simulación, como son la amplitud del campo magnético, eléctrico o el campo de velocidades lo cual ayudará a confirmar o descartar las hipótesis planteada de como la resistividad afecta en el crecimiento de las inestabilidades.

El estudio de las inestabilidades KHI es relevante para entender fenómenos que se dan en sistemas astrofísicos como chorros relativistas, magnetosferas ó estrellas de neutrones. A través del análisis de inestabilidades como la KHI, se puede obtener información valiosa sobre el comportamiento de estos sistemas extremos y los procesos físicos que ocurren en ellos.